%%%%%%%%%%%%%%%%%%%%%%%%%%%%%%%%%%%%%%%%%%%%%%%%%%%
%
%  New template code for TAMU Theses and Dissertations starting Spring 2021.
%
%
%  Author: Thesis Office
%
%  Last Updated: 1/13/2021
%
%%%%%%%%%%%%%%%%%%%%%%%%%%%%%%%%%%%%%%%%%%%%%%%%%%%
%%%%%%%%%%%%%%%%%%%%%%%%%%%%%%%%%%%%%%%%%%%%%%%%%%%%%%%%%%%%%%%%%%%%%
%%                           ABSTRACT
%%%%%%%%%%%%%%%%%%%%%%%%%%%%%%%%%%%%%%%%%%%%%%%%%%%%%%%%%%%%%%%%%%%%%

\chapter*{ABSTRACT}
\addcontentsline{toc}{chapter}{ABSTRACT} % Needs to be set to part, so the TOC doesn't add 'CHAPTER ' prefix in the TOC.

\pagestyle{plain} % No headers, just page numbers
\pagenumbering{roman} % Roman numerals
\setcounter{page}{2}

\indent
%This dissertation develops precision inference methods and optical instrumentation that reduce hardware complexity in spectroscopic systems while maintaining accuracy and sensitivity. Interferometric techniques are first introduced to characterize laser wavelength and detuning with high precision and are validated through atomic hyperfine spectroscopy of rubidium and xenon. These methods are then used to validate calibration-free spectroscopic approaches in which wavelength information is inferred directly from spectral structure, enabling the elimination of dedicated wavelength measurement hardware. Cavity-enhanced spectroscopic systems are developed to achieve high sensitivity in compact geometries, with deterministic spectral fitting and uncertainty estimation. Applications include trace-gas detection and the empirical determination of molecular absorption cross sections, along with the identification of systematic effects arising from surface–molecule interactions. Finally, novel instrumentation enabling fast, broadband spectroscopy is presented, further demonstrating how precision inference can enable simple, robust, and compact optical spectrometers.
Optical spectroscopy offers highly selective and sensitive molecular detection, yet its use in compact and practical sensing platforms is often limited by the complexity and calibration requirements associated with wavelength control and measurement. These challenges are especially pronounced in breath analysis, where weak absorption features, spectral congestion, and the need for rapid, quantitative measurements place stringent demands on both sensitivity and wavelength stability. This dissertation develops and validates complementary spectroscopic techniques that reduce hardware complexity while preserving quantitative accuracy through inference-based and interferometric approaches.

Precision methods for laser wavelength and optical detuning measurement are first established using a Michelson interferometer and validated through atomic hyperfine spectroscopy. Building on this foundation, calibration-free spectroscopic inference techniques are developed that simultaneously recover wavelength detuning and molecular species concentrations directly from measured spectral structure, eliminating the need for independent wavelength measurement instrumentation in compatible systems. These approaches are implemented within cavity-enhanced spectroscopic platforms and applied to measurements of human breath samples. The work is further extended to broadband cavity-enhanced spectroscopy using optical frequency combs and Vernier coupling, enabling rapid spectral acquisition with relaxed detector resolution requirements. For this platform, a novel spatially resolved spectrometer detector employing a single photodiode is developed. Together, these contributions demonstrate a unified strategy for realizing compact, flexible, and high-throughput spectroscopic systems suitable for quantitative chemical sensing and breath analysis.



\pagebreak{}
